% ============================================================================
% COMPREHENSIVE PREAMBLE FOR ERDOS PROBLEM 915 THESIS
% ============================================================================

\documentclass[kul, 11pt, a4paper, oneside]{kulakreport}
\usetikzlibrary{calc, arrows.meta, shapes.geometric, decorations.markings}

% KU Leuven metadata
\faculty{Faculty of Science}
\group{Department of Mathematics}
\subtitle{Erd\H{o}s Problem 915, from Proof to Search}
\institute{Master of Mathematics}
\emailaddress{louis.vandenbruwaene@student.kuleuven.be}
\address{
   \textbf{KU Leuven} \\
   \textbf{Faculty of Science} \\
   Celestijnenlaan 200B, 3001 Leuven \\
   Belgium
}

% ============================================================================
% BASIC PACKAGES
% ============================================================================
\usepackage[utf8]{inputenc}
\usepackage[T1]{fontenc}
\usepackage[scaled=0.95]{helvet}
\renewcommand{\familydefault}{\sfdefault}
\usepackage[absolute,overlay]{textpos}
% lmodern not available; disable microtype expansion to avoid font errors
\usepackage[patch=none, expansion=false]{microtype}

% ============================================================================
% PAGE LAYOUT
% ============================================================================
\usepackage[
    a4paper,
    inner=25mm,
    outer=25mm,
    top=30mm,
    bottom=30mm,
    bindingoffset=5mm
]{geometry}

\usepackage{setspace}
\onehalfspacing

% Improve line breaking to reduce overfull hbox warnings
\tolerance=1414
\emergencystretch=1.5em
\hfuzz=3pt
\vfuzz=3pt

% ============================================================================
% MATHEMATICS
% ============================================================================
\usepackage{amsmath, amssymb, amsthm}
\usepackage{mathtools}
\usepackage{bm}

% ============================================================================
% GRAPHICS AND COLOR (tikz, graphicx, xcolor already loaded by kulakreport)
% ============================================================================
\usetikzlibrary{positioning, arrows.meta, calc, shapes.geometric,
                decorations.pathmorphing, backgrounds, fit,
                patterns, shadows, matrix}

% ============================================================================
% FIGURES AND TABLES
% ============================================================================
\usepackage{float}
% rotating: the twelve-panel bounds grid needs a full sideways page, since at
% text width its per-panel labels fall under 2pt.
\usepackage{rotating}
\usepackage{caption}
\usepackage{subcaption}
\usepackage{booktabs}
\usepackage{array}
\usepackage{multirow}

% Uniform caption look for every plot, figure, and table: a bold blue label, a
% slightly smaller justified body set in from the margins, and a touch of space
% above so a caption never crowds the artwork it describes.
\captionsetup{
    font=small,
    labelfont={bf,color=KULblauw3b},
    textfont={color=black!85},
    labelsep=period,
    justification=justified,
    singlelinecheck=true,
    margin=14pt,
    skip=8pt,
}
\captionsetup[sub]{
    font=footnotesize,
    labelfont={bf,color=KULblauw3b},
    % Ragged right, not justified: sub-captions often sit in narrow boxes
    % (half-width panels), where justification stretches short lines into
    % wide word gaps around unbreakable math.
    justification=raggedright,
    margin=4pt,
    skip=5pt,
}

% ============================================================================
% LISTS
% ============================================================================
\usepackage{enumitem}

% Custom bullet styles using small graph icons
\newcommand{\prettygrotzsch}{%
	\tikz[baseline=-0.6ex, scale=0.12]{
		\definecolor{nodecolor}{RGB}{50, 120, 180}
		\foreach \i in {0,...,4} {
			\coordinate (inner\i) at ({1.2*cos(72*\i+18)},{1.2*sin(72*\i+18)});
		}
		\coordinate (center) at (0,0);
		\draw[nodecolor, line width=0.5pt]
		(center) -- (inner0) (center) -- (inner1) (center) -- (inner2)
		(center) -- (inner3) (center) -- (inner4)
		(inner0) -- (inner2) (inner0) -- (inner3) (inner1) -- (inner3)
		(inner1) -- (inner4) (inner2) -- (inner4);
		\foreach \i in {0,...,4} { \fill[black] (inner\i) circle (2.5pt); }
		\fill[black] (center) circle (2.5pt);
	}
}

\newcommand{\prettypetersen}{%
	\tikz[baseline=-0.6ex, scale=0.15]{
		\definecolor{petersennode}{RGB}{20, 140, 140}
		\definecolor{petersenedge}{RGB}{100, 100, 100}

		\def\R{1.2}
		\def\r{0.8}

		\foreach \i in {0,...,4} {
			\coordinate (outer\i) at ({90+72*\i}:\R);
			\coordinate (inner\i) at ({90+72*\i}:\r);
		}

		\foreach \i in {0,...,4} {
			\fill[petersennode] (outer\i) circle (4pt);
			\fill[petersennode] (inner\i) circle (4pt);
		}

		\draw[petersenedge, line width=0.7pt]
		(outer0) -- (outer1) -- (outer2) -- (outer3) -- (outer4) -- (outer0)
		(outer0) -- (inner0)
		(outer1) -- (inner1)
		(outer2) -- (inner2)
		(outer3) -- (inner3)
		(outer4) -- (inner4)
		(inner0) -- (inner2) -- (inner4) -- (inner1) -- (inner3) -- (inner0);

		\foreach \i in {0,...,4} {
			\fill[black, opacity=1] (outer\i) circle (3pt);
			\fill[black, opacity=1] (inner\i) circle (3pt);
		}
	}
}

\setlist[itemize]{label=$\prettygrotzsch$, leftmargin=1.5em}
\setlist[enumerate]{label=\protect\prettypetersen, leftmargin=2em}

% ============================================================================
% HYPERLINKS AND REFERENCES
% ============================================================================
\usepackage[noabbrev, capitalize, nameinlink]{cleveref}

% ============================================================================
% BIBLIOGRAPHY
% ============================================================================
\usepackage[numbers, sort&compress]{natbib}

% ============================================================================
% HEADER/FOOTER
% ============================================================================
% The page footer ("section mark | page number") is defined by the kulakreport
% class with oneside-friendly specifiers; nothing to add here.

% ============================================================================
% CHAPTER STYLING
% ============================================================================
\usepackage{titlesec}
% Chapter: keep the large black title, but brand the "Chapter N" label in the
% KU Leuven house blue so the heading reads as part of the institutional style.
\titleformat{\chapter}[display]
  {\normalfont\huge\bfseries}{\color{KULblauw1}\chaptertitlename\ \thechapter}{20pt}{\Huge\raggedright}
\titlespacing*{\chapter}{0pt}{-30pt}{40pt}

% Section headings: black bold title with a blue section number and a thin blue
% accent rule beneath, matching the chapter branding and easing navigation in
% the long chapters. Subsection and below stay quieter, with only a blue number.
\titleformat{\section}
  {\normalfont\Large\bfseries}{\color{KULblauw1}\thesection}{0.7em}{}
\titlespacing*{\section}{0pt}{1.8ex plus .3ex minus .2ex}{0.9ex}

\titleformat{\subsection}
  {\normalfont\large\bfseries}{\color{KULblauw1}\thesubsection}{0.6em}{}
\titlespacing*{\subsection}{0pt}{1.4ex plus .2ex minus .2ex}{0.6ex}

\titleformat{\subsubsection}
  {\normalfont\normalsize\bfseries}{\color{KULblauw1}\thesubsubsection}{0.5em}{}
\titlespacing*{\subsubsection}{0pt}{1.1ex plus .2ex minus .2ex}{0.5ex}

\titleformat{\paragraph}[runin]
  {\normalfont\normalsize\bfseries\color{KULblauw3b}}{}{0pt}{}
\titlespacing*{\paragraph}{0pt}{1ex plus .2ex minus .1ex}{0.8em}

% Part openers carry a framing paragraph rather than sitting as an empty page.
% The "top" class starts a new page with the banner at the top, then lets the
% following text flow on the same page.
\titleclass{\part}{top}
\titleformat{\part}[display]
  {\normalfont\Huge\bfseries\color{KULblauw3b}}{\partname\ \thepart}{0.4em}{\Huge}
\titlespacing*{\part}{0pt}{0pt}{1.5em}

% ============================================================================
% THEOREM ENVIRONMENTS (using aliascnt for proper cleveref support)
% ============================================================================
\usepackage{aliascnt}

% Base theorem
\theoremstyle{plain}
\newtheorem{theorem}{Theorem}[chapter]

% Theorem-like (italic body) - with aliascnt for cleveref
\newaliascnt{proposition}{theorem}
\newtheorem{proposition}[proposition]{Proposition}
\aliascntresetthe{proposition}

\newaliascnt{lemma}{theorem}
\newtheorem{lemma}[lemma]{Lemma}
\aliascntresetthe{lemma}

\newaliascnt{corollary}{theorem}
\newtheorem{corollary}[corollary]{Corollary}
\aliascntresetthe{corollary}

\newaliascnt{conjecture}{theorem}
\newtheorem{conjecture}[conjecture]{Conjecture}
\aliascntresetthe{conjecture}

\newaliascnt{claim}{theorem}
\newtheorem{claim}[claim]{Claim}
\aliascntresetthe{claim}

% Definition-like (upright body)
\theoremstyle{definition}
\newaliascnt{definition}{theorem}
\newtheorem{definition}[definition]{Definition}
\aliascntresetthe{definition}

\newaliascnt{example}{theorem}
\newtheorem{example}[example]{Example}
\aliascntresetthe{example}

\newaliascnt{construction}{theorem}
\newtheorem{construction}[construction]{Construction}
\aliascntresetthe{construction}

% Remark-like
\theoremstyle{remark}
\newaliascnt{remark}{theorem}
\newtheorem{remark}[remark]{Remark}
\aliascntresetthe{remark}

\newaliascnt{notation}{theorem}
\newtheorem{notation}[notation]{Notation}
\aliascntresetthe{notation}

\newaliascnt{ques}{theorem}
\newtheorem{ques}[ques]{Question}
\aliascntresetthe{ques}

% Unnumbered versions
\newtheorem*{theorem*}{Theorem}
\newtheorem*{remark*}{Remark}

% Custom proof environment for claims
\newenvironment{claimproof}[1][\proofname]{\begin{proof}[#1]\renewcommand{\qedsymbol}{$\diamondsuit$}}{\end{proof}}

% Cleveref settings (must be after theorem environment definitions)
\crefname{theorem}{theorem}{theorems}
\Crefname{theorem}{Theorem}{Theorems}
\crefname{proposition}{proposition}{propositions}
\Crefname{proposition}{Proposition}{Propositions}
\crefname{lemma}{lemma}{lemmas}
\Crefname{lemma}{Lemma}{Lemmas}
\crefname{corollary}{corollary}{corollaries}
\Crefname{corollary}{Corollary}{Corollaries}
\crefname{conjecture}{conjecture}{conjectures}
\Crefname{conjecture}{Conjecture}{Conjectures}
\crefname{definition}{definition}{definitions}
\Crefname{definition}{Definition}{Definitions}
\crefname{example}{example}{examples}
\Crefname{example}{Example}{Examples}
\crefname{remark}{remark}{remarks}
\Crefname{remark}{Remark}{Remarks}
\crefname{construction}{construction}{constructions}
\Crefname{construction}{Construction}{Constructions}
\crefname{ques}{question}{questions}
\Crefname{ques}{Question}{Questions}
\crefname{chapter}{Chapter}{Chapters}
\crefname{appendix}{Appendix}{Appendices}
\Crefname{appendix}{Appendix}{Appendices}
\crefname{section}{Section}{Sections}
\crefname{part}{Part}{Parts}
\crefname{figure}{Figure}{Figures}
\crefname{table}{Table}{Tables}

% ============================================================================
% CUSTOM MATHEMATICAL COMMANDS
% ============================================================================

% Floor and ceiling
\newcommand*{\floor}[1]{\left\lfloor #1 \right\rfloor}
\newcommand*{\ceil}[1]{\left\lceil #1 \right\rceil}
\newcommand*{\floorfrac}[2]{\floor{\frac{#1}{#2}}}
\newcommand*{\ceilfrac}[2]{\ceil{\frac{#1}{#2}}}

% Absolute value and norms
\newcommand*{\abs}[1]{\left\lvert #1 \right\rvert}
\newcommand*{\norm}[1]{\left\lVert #1 \right\rVert}

% Graph theory notation
\newcommand{\lec}{\lambda}                    % Local edge-connectivity
\newcommand{\kap}{\kappa}                     % Local vertex-connectivity
\newcommand{\multigraph}{\mathcal{M}}         % Multigraph
\newcommand{\HH}{\mathcal{H}}                 % Hypergraph
\DeclareMathOperator{\dist}{dist}             % Distance

% Probability notation
\DeclareMathOperator{\Var}{Var}               % Variance

% Common sets
\newcommand{\N}{\mathbb{N}}
\newcommand{\Z}{\mathbb{Z}}
\newcommand{\R}{\mathbb{R}}

% ============================================================================
% TIKZ STYLES - Unified Colors and Styles (from shared/)
% ============================================================================

% Import shared color definitions and TikZ styles
% ============================================================================
% UNIFIED COLOR DEFINITIONS FOR ERDOS PROBLEM 915 THESIS
% ============================================================================
% This file provides consistent colors across thesis, progress report, and
% presentation. All colors are based on KU Leuven institutional branding.
% ============================================================================

% --- KU Leuven Primary Colors ---
% Note: In beamer (kulakbeamer), KULblauw1/KULblauw3a/KULblauw3b are predefined
% For thesis documents, we define them here
\makeatletter
\@ifclassloaded{kulakbeamer}{%
    % Colors already defined by beamer class
}{%
    % Define KUL colors for non-beamer documents
    \definecolor{KULblauw1}{RGB}{29,141,176}     % Main blue
    \definecolor{KULblauw3a}{RGB}{82,189,236}    % Accent light blue
    \definecolor{KULblauw3b}{RGB}{30,110,135}    % Accent dark blue
}
\makeatother

% --- Graph Vertex Colors ---
% Primary vertex color uses KUL blue
\definecolor{vertexblue}{RGB}{29,141,176}        % KULblauw1
\definecolor{vertexred}{RGB}{200,80,80}          % Matched with presentation
\definecolor{vertexpurple}{RGB}{140,80,160}      % Matched with presentation
\definecolor{vertexgreen}{RGB}{60,160,80}        % Matched with presentation
\definecolor{vertexorange}{RGB}{220,140,40}      % Matched with presentation

% --- Aliases for presentation compatibility ---
\colorlet{graphred}{vertexred}
\colorlet{graphpurple}{vertexpurple}
\colorlet{graphgreen}{vertexgreen}
\colorlet{graphorange}{vertexorange}

% ============================================================================
% UNIFIED TIKZ STYLES FOR ERDOS PROBLEM 915 THESIS
% ============================================================================
% This file provides consistent TikZ graph styles across thesis, progress
% report, and presentation. Requires colors.tex to be loaded first.
% ============================================================================

% --- Vertex Styles ---
\tikzset{
    % Basic vertex style (KUL blue)
    vertex/.style={
        circle, draw=KULblauw1, fill=KULblauw1!15,
        minimum size=20pt, inner sep=0pt,
        font=\small\bfseries, line width=1.5pt
    },
    % Small vertex for complex diagrams
    smallvertex/.style={
        circle, draw=KULblauw1, fill=KULblauw1!15,
        minimum size=16pt, inner sep=0pt,
        font=\scriptsize\bfseries, line width=1.2pt
    },
    % Colored variants - semantic naming
    leftvertex/.style={vertex, draw=KULblauw1, fill=KULblauw1!15},
    rightvertex/.style={vertex, draw=vertexred, fill=vertexred!15},
    sharedvertex/.style={vertex, draw=vertexpurple, fill=vertexpurple!20},
    tailvertex/.style={vertex, draw=vertexgreen, fill=vertexgreen!15},
    hubvertex/.style={vertex, draw=vertexorange, fill=vertexorange!20},
    % Small colored variants
    smallleftvertex/.style={smallvertex, draw=KULblauw1, fill=KULblauw1!15},
    smallrightvertex/.style={smallvertex, draw=vertexred, fill=vertexred!15},
    smallsharedvertex/.style={smallvertex, draw=vertexpurple, fill=vertexpurple!20},
    smalltailvertex/.style={smallvertex, draw=vertexgreen, fill=vertexgreen!15},
    % Presentation-style aliases (gnode = graph node)
    gnode/.style={vertex},
    bluenode/.style={leftvertex},
    rednode/.style={rightvertex},
    purplenode/.style={sharedvertex},
    greennode/.style={tailvertex},
    orangenode/.style={hubvertex},
}

% --- Edge Styles ---
\tikzset{
    % Basic edge styles
    edge/.style={line width=1.8pt, KULblauw1!70},
    thickedge/.style={line width=2.5pt},
    % Colored edge variants - semantic naming
    leftedge/.style={line width=1.8pt, KULblauw1},
    rightedge/.style={line width=1.8pt, vertexred},
    sharededge/.style={line width=1.8pt, vertexpurple!80},
    tailedge/.style={line width=1.8pt, vertexgreen!80!black},
    % Tree edges
    treeedge/.style={line width=1.8pt, KULblauw3b},
    lefttreeedge/.style={line width=1.8pt, KULblauw1},
    righttreeedge/.style={line width=1.8pt, vertexred},
    % Special edge styles
    missingedge/.style={line width=1.5pt, dashed, gray},
    multiedge/.style={line width=1.5pt, draw=KULblauw3b},
    hyperedge/.style={draw=vertexred, fill=vertexred!10, rounded corners=8pt, line width=1.5pt},
    % Presentation-style aliases
    gedge/.style={line width=2pt, KULblauw1!70},
    blueedge/.style={line width=2pt, KULblauw1},
    rededge/.style={line width=2pt, vertexred},
    purpleedge/.style={line width=2pt, vertexpurple!80},
    greenedge/.style={line width=2pt, vertexgreen!80!black},
    tedge/.style={treeedge},
}

% --- Directed arc styles (Stealth arrowheads, thesis blue) ---
% A global `arc` is safe: figures that set their own `arc/.style` inline
% (e.g. the Chapter 1 hub figures) shadow this within their own picture.
\tikzset{
    arc/.style={-{Stealth[length=2.2mm]}, line width=1.1pt, KULblauw1},
    faintarc/.style={-{Stealth[length=1.8mm]}, line width=0.7pt, KULblauw1!35},
    redarc/.style={-{Stealth[length=2.2mm]}, line width=1.1pt, dashed, vertexred},
    forbiddenedge/.style={line width=1.3pt, dashed, vertexred},
}

% --- Hyperedge node: a hollow square standing for a hyperedge (e.g. in an
%     incidence graph), distinct from the `hyperedge` region style above ---
\tikzset{
    hnode/.style={
        rectangle, draw=KULblauw3b, fill=white,
        minimum size=13pt, inner sep=1pt, line width=1pt,
        font=\scriptsize\bfseries
    },
}

% --- Label and Weight Styles ---
\tikzset{
    treeweight/.style={
        circle, fill=white, inner sep=1.5pt,
        font=\scriptsize\bfseries
    },
    tweight/.style={
        circle, fill=white, inner sep=2pt,
        font=\small\bfseries
    },
    edgelabel/.style={
        font=\scriptsize, fill=white, inner sep=1pt
    }
}


% Two further metro-line colours, defined thesis-local (the shared file stops at
% metroF) so the larger hyperedge figures in Chapter 2 can give every one of
% their eight lines its own distinct hue.
\definecolor{metroG}{RGB}{200,70,140}   % magenta line
\definecolor{metroH}{RGB}{30,130,120}   % teal line

% Thesis-local graph polish. These refinements are applied here, after the
% shared file, so the slides and progress report keep their own look untouched.
% A faint drop shadow lifts every vertex off the page for a cleaner, more
% three-dimensional read, and rounded line caps soften the stroke ends of the
% edges. Both are deliberately subtle so dense incidence diagrams stay legible.
\usetikzlibrary{shadows}
\tikzset{
    vertex/.append style={
        drop shadow={shadow xshift=0.5pt, shadow yshift=-0.6pt,
                     fill=black, opacity=0.20}},
    smallvertex/.append style={
        drop shadow={shadow xshift=0.4pt, shadow yshift=-0.5pt,
                     fill=black, opacity=0.18}},
    hnode/.append style={
        drop shadow={shadow xshift=0.4pt, shadow yshift=-0.5pt,
                     fill=black, opacity=0.16}},
    edge/.append style={line cap=round},
    thickedge/.append style={line cap=round},
    leftedge/.append style={line cap=round},
    rightedge/.append style={line cap=round},
    sharededge/.append style={line cap=round},
    tailedge/.append style={line cap=round},
    treeedge/.append style={line cap=round},
    lefttreeedge/.append style={line cap=round},
    righttreeedge/.append style={line cap=round},
    multiedge/.append style={line cap=round},
    gedge/.append style={line cap=round},
    blueedge/.append style={line cap=round},
    rededge/.append style={line cap=round},
    purpleedge/.append style={line cap=round},
    greenedge/.append style={line cap=round},
}

% ============================================================================
% CANONICAL THESIS GRAPH VOCABULARY
% ============================================================================
% One visual language for every graph in the thesis, body and appendix alike, so
% the figures read as a family instead of each hand-rolling its own look. It is
% grouped by the five kinds of drawing the thesis does: SIMPLE graphs, MULTI
% graphs, DIRECTED graphs, CURVED connections, and the METRO-map hyperedges.
% A figure may override the size for a dense diagram (e.g. gvertex, minimum
% size=4.5mm), but the colour, stroke, fill, font, and shadow stay fixed so the
% family stays coherent. Curves are never removed here: every bend lives in a
% figure's own \draw ... to[bend ...], and the styles below fix only how a line
% or arc looks, not whether it bends.
\tikzset{
  % ---- SIMPLE: the default blue vertex and the plain undirected edge ----
  gvertex/.style={circle, draw=KULblauw1, fill=KULblauw1!15, line width=1pt,
                  minimum size=6mm, inner sep=0pt, font=\small,
                  drop shadow={shadow xshift=0.4pt, shadow yshift=-0.5pt,
                               fill=black, opacity=0.18}},
  gline/.style ={line width=1.3pt, KULblauw1!75, line cap=round},
  gfaint/.style={line width=0.9pt, KULblauw1!40, line cap=round},
  % semantic vertex colours: same size and stroke, the colour marks the role
  gvhub/.style ={gvertex, draw=vertexorange, fill=vertexorange!18},
  gvsink/.style={gvertex, draw=vertexred,    fill=vertexred!15},
  gvcut/.style ={gvertex, draw=vertexpurple, fill=vertexpurple!20},
  gvnew/.style ={gvertex, draw=vertexorange, fill=vertexorange!22},
  % ---- MULTI: parallel edges are simple edges drawn with a bend (see CURVED) ---
  gmulti/.style={gline},
  % ---- DIRECTED: arcs carry a Stealth head; faint background and cut variants --
  gdir/.style     ={-{Stealth[length=2.2mm]}, line width=1.1pt, KULblauw1!75, line cap=round},
  gdirfaint/.style={-{Stealth[length=1.8mm]}, line width=0.8pt, KULblauw1!40, line cap=round},
  gcut/.style     ={-{Stealth[length=2.2mm]}, line width=1.2pt, densely dashed, vertexred},
  gcutedge/.style ={line width=1.2pt, densely dashed, vertexred},
  % highlighted routes: three independent routes always read orange, green, purple
  grA/.style={line width=2.6pt, vertexorange,        line cap=round},
  grB/.style={line width=2.6pt, vertexgreen!70!black, line cap=round},
  grC/.style={line width=2.6pt, vertexpurple,         line cap=round},
  grAd/.style={grA, -{Stealth[length=2.8mm]}},
  grBd/.style={grB, -{Stealth[length=2.8mm]}},
  grCd/.style={grC, -{Stealth[length=2.8mm]}},
  % ---- CURVED: the standard gentle bend for a parallel pair or a 2-cycle ----
  gcurve/.style ={bend left=14},
  gcurveR/.style={bend right=14},
  % ---- METRO: hyperedges as metro lines. The metro, metrostop, and metroA..H
  %      styles live in shared/tikz-styles.tex and are reused here unchanged, so
  %      the hyperedge look is already shared across every figure that uses it. --
  %
  % Re-point the shared graph names the thesis body already uses (vertex,
  % smallvertex, edge) onto the canonical look, so Chapter 2, Chapter 4, and any
  % figure using them cohere with the rest without per-figure edits. This is
  % thesis-local: the slides load shared/tikz-styles.tex directly and keep their
  % own sizes. smallvertex holds its size and small font so multi-character split
  % labels (v^in, v^out) still fit.
  vertex/.style={gvertex},
  smallvertex/.style={gvertex, minimum size=16pt, font=\scriptsize},
  edge/.style={gline},
  arc/.style={gdir},
}

% Link appearance. hyperref (loaded by the class) otherwise draws a coloured box
% around every cross-reference, citation, and table-of-contents entry, which is
% distracting in print. Switch to coloured text with no boxes, using the dark
% KU Leuven blue of the chapter headings so links read cleanly on paper and on
% screen. Placed here because it must come after the KUL colours are defined.
\hypersetup{
    colorlinks=true,
    linkcolor=KULblauw3b,
    citecolor=KULblauw3b,
    urlcolor=KULblauw3b,
    linktoc=all,
}

% ============================================================================
% DOCUMENT METADATA
% ============================================================================
\title{Connectivity Bounds in Graphs and Beyond}
\author{Louis Vandenbruwaene}
\date{Academic Year 2026--2027}

% PDF document properties, so the file identifies itself in a reader's title bar,
% in a library catalogue, and in any tool that reads the metadata rather than the
% cover page. Without these the fields are blank.
\hypersetup{
    pdftitle={Connectivity Bounds in Graphs and Beyond},
    pdfauthor={Louis Vandenbruwaene},
    pdfsubject={Extremal graph theory: Erdos Problem 915 across twelve structural variants},
    pdfkeywords={extremal graph theory, edge-disjoint paths, vertex-disjoint paths,
                 local connectivity, Menger, Gomory-Hu, hypergraphs, digraphs,
                 Erdos Problem 915, mixed-integer linear programming, tabu search},
    pdfcreator={pdfTeX},
}

% The submitted snapshot. revision.tex is written by ./record_revision.sh, which
% also tags the same commit, so the printed hash and the tag always agree. When
% the file is absent (a working build between revisions) the appendix says so
% rather than printing a stale hash.
\newcommand{\thesistag}{submitted}
\newcommand{\thesiscommit}{not recorded in this build}
\IfFileExists{revision.tex}{\input{revision.tex}}{}

% ============================================================================
% OFFICIAL FACULTY OF SCIENCE COVER
% ============================================================================
\definecolor{bluetitle}{RGB}{29,141,176}
\definecolor{blueaff}{RGB}{0,0,128}
\definecolor{blueline}{RGB}{82,189,236}
\setlength{\TPHorizModule}{1mm}
\setlength{\TPVertModule}{1mm}
% The faculty-cover coordinates below were written for the standard 1in TeX
% text origin, but textpos is loaded with [absolute] (origin = physical page
% corner). Without resetting the origin, the negative-coordinate header row
% (banner, logo, address) is pushed off the top-left edge of the page. Pin the
% origin back to the 1in reference so both covers land where they were designed.
\textblockorigin{1in}{1in}

\newcommand{\makefacultycover}{%
\thispagestyle{empty}%
\begingroup
\hfuzz=40pt
\sffamily
\begin{textblock}{191}(-24,-11)
\colorbox{bluetitle}{\hspace{139mm}\ \parbox[c][18truemm]{52mm}{\textcolor{white}{FACULTY OF SCIENCE}}}
\end{textblock}
\begin{textblock}{70}(-18,-19)
\includegraphics*[height=19.8truemm]{LogoKULeuven}
\end{textblock}
\begin{textblock}{160}(8,63)
\flushleft
\fontsize{34}{38}\selectfont \textcolor{bluetitle}{\thetitle}\\[2mm]
\fontsize{18}{22}\selectfont \thesubtitle
\end{textblock}
\begin{textblock}{160}(8,153)
\flushright
\fontsize{14}{16}\selectfont \textbf{\theauthor}
\end{textblock}
\begin{textblock}{80}(8,187)
\flushleft
Promotor: Prof.\ Jan Goedgebeur\\[-2pt]
Supervisor: Prof.\ Stijn Cambie\\[-2pt]
\textcolor{blueaff}{KU Leuven}\\
\end{textblock}
\begin{textblock}{160}(8,191)
\flushright
Thesis presented in\\[4.5pt]
fulfillment of the requirements\\[4.5pt]
for the degree of Master of Science\\[4.5pt]
in Mathematics\\
\end{textblock}
\begin{textblock}{160}(8,232)
\flushright
Academic year 2026--2027
\end{textblock}
\begin{textblock}{191}(-24,248)
{\color{blueline}\rule{550pt}{5.5pt}}
\end{textblock}
\mbox{}
\endgroup
\clearpage
}

\newcommand{\makefacultybackcover}{%
\clearpage
\ifodd\value{page}\else\mbox{}\thispagestyle{empty}\clearpage\fi
\thispagestyle{empty}%
\begingroup
\hfuzz=40pt
\sffamily
\begin{textblock}{191}(113,-11)
{\color{blueline}\rule{160pt}{5.5pt}}
\end{textblock}
\begin{textblock}{191}(168,-11)
{\color{blueline}\rule{5.5pt}{59pt}}
\end{textblock}
\begin{textblock}{183}(-24,-11)
\flushright
\fontsize{7}{7.5}\selectfont
\textbf{DEPARTMENT OF MATHEMATICS}\\
Celestijnenlaan 200B\\
3001 Leuven, Belgium\\
wis.kuleuven.be\\
\end{textblock}
\begin{textblock}{191}(154,-7)
\includegraphics*[height=16.5truemm]{sedes}
\end{textblock}
\begin{textblock}{191}(-20,235)
{\color{bluetitle}\rule{544pt}{55pt}}
\end{textblock}
\mbox{}
\endgroup
}

% ============================================================================
% COPYRIGHT AND DISCLAIMER PAGE
% ----------------------------------------------------------------------------
% Required by the Faculty of Science layout rules: the first page after the
% cover carries the copyright notice, and the open-access disclaimer marking
% the text as student work must follow the title page.
% https://wet.kuleuven.be/english/mastersthesis/form
% ============================================================================
\newcommand{\makecopyrightpage}{%
\thispagestyle{empty}%
\begingroup
\sffamily
\null\vfill
\begingroup
\footnotesize
\noindent\textbf{Copyright information}\\[2pt]
Student paper as part of an academic education and examination. No correction
was made to the paper after examination.

\bigskip
\noindent\textcopyright\ Copyright by KU Leuven

\medskip
\noindent Without prior written permission from both the supervisor(s) and the
author(s), copying, reproducing, using, or realizing this publication or parts
thereof is prohibited. For requests or information regarding the copying and/or
use and/or realization of parts of this publication, please contact KU Leuven,
Faculty of Science, Celestijnenlaan 200H box 2100, 3001 Leuven (Heverlee),
Telephone +32 16 32 14 01.

\medskip
\noindent Prior written permission from the supervisor(s) is also required for
using the (original) methods, products, circuits, and programs described in this
work for industrial or commercial purposes, and for submitting this publication
for scientific awards or contests.
\par
\endgroup
\endgroup
\clearpage
}

% ============================================================================
% MISCELLANEOUS
% ============================================================================

% Paragraph spacing
\setlength{\parskip}{0.25em}
\setlength{\parindent}{1.5em}
\setcounter{tocdepth}{2}

% Better theorem spacing
\makeatletter
\def\thm@space@setup{%
    \thm@preskip=0.75em plus 0.25em minus 0.2em
    \thm@postskip=0.75em plus 0.25em minus 0.2em
}
\makeatother

% Let long display chains break across a page rather than pushing a whole
% block to the next one, which otherwise leaves half-empty pages.
\allowdisplaybreaks[2]

% Float separation: the defaults leave more air around figures than this
% document's dense pages need.
\setlength{\floatsep}{10pt plus 3pt minus 2pt}
\setlength{\textfloatsep}{12pt plus 3pt minus 2pt}
\setlength{\intextsep}{10pt plus 3pt minus 2pt}

% Epigraph for chapter quotes
\usepackage{epigraph}
\setlength{\epigraphwidth}{0.7\textwidth}

% For including appendices
\usepackage{appendix}

% ============================================================================
% ALGORITHM BOX ENVIRONMENT
% ============================================================================
\usepackage{tcolorbox}
\tcbuselibrary{skins, breakable}

\newtcolorbox{algorithmbox}[1][]{%
    enhanced,
    breakable,
    colback=KULblauw1!5,
    colframe=KULblauw1!70,
    fonttitle=\bfseries,
    left=2mm, right=2mm, top=1mm, bottom=1mm,
    before upper={\parindent0pt\ttfamily\small},
    #1
}

% ============================================================================
% PSEUDOCODE (algorithm2e) AND CODE LISTINGS (listings)
% ============================================================================
% These power the computational-centerpiece chapters: pseudocode for the
% oracle / annealing / certifier, and Python listings in the code appendix.
\usepackage[ruled, vlined, linesnumbered]{algorithm2e}
% On-brand pseudocode: a blue "Algorithm N" label, keywords in the house blue,
% comments in the accent blue, and grey line numbers, so the algorithm floats
% read in the same palette as the code cards and listings. The ruled (KU
% Leuven-safe) layout is kept; only the typography is recoloured. The style
% setters take one-argument command names, so we wrap the colours in macros.
\newcommand{\algokwsty}[1]{\textcolor{KULblauw1}{\textbf{#1}}}
\newcommand{\algofuncsty}[1]{\textcolor{KULblauw1}{\textbf{#1}}}
\newcommand{\algocommentsty}[1]{\textcolor{KULblauw3b}{\textit{#1}}}
\SetAlCapNameFnt{\color{KULblauw3b}\bfseries}
\SetAlCapFnt{\small}
\SetKwSty{algokwsty}
\SetFuncSty{algofuncsty}
\SetCommentSty{algocommentsty}
\usepackage{listings}

% A calm, on-brand style for the Python listings of the unified program.
\lstdefinestyle{pythonkul}{%
    language=Python,
    basicstyle=\ttfamily\footnotesize,
    keywordstyle=\color{KULblauw1}\bfseries,
    commentstyle=\color{KULblauw3b}\itshape,
    stringstyle=\color{vertexgreen},
    numberstyle=\ttfamily\scriptsize\color{gray},
    backgroundcolor=\color{KULblauw1!4},
    numbers=left,
    numbersep=8pt,
    breaklines=true,
    breakatwhitespace=false,
    postbreak=\mbox{\textcolor{KULblauw3a}{$\hookrightarrow$}\space},
    showstringspaces=false,
    frame=leftline,
    framerule=1pt,
    rulecolor=\color{KULblauw3a},
    xleftmargin=12pt,
    tabsize=4,
    captionpos=b,
    % The program's figure labels contain a few non-ASCII characters (the Greek
    % sigma and an arrow). When listed verbatim under utf8 inputenc these would
    % be "characters not set up" and abort the build, so map each to its math
    % rendering. A handful of related symbols are included for robustness.
    inputencoding=utf8,
    extendedchars=true,
    literate=%
        {σ}{{$\sigma$}}1 {λ}{{$\lambda$}}1 {κ}{{$\kappa$}}1 {μ}{{$\mu$}}1
        {α}{{$\alpha$}}1 {Λ}{{$\Lambda$}}1 {→}{{$\rightarrow$}}1
        {≤}{{$\leq$}}1 {≥}{{$\geq$}}1 {≠}{{$\neq$}}1 {×}{{$\times$}}1
        {∈}{{$\in$}}1 {∞}{{$\infty$}}1 {·}{{$\cdot$}}1
        {ő}{{\H{o}}}1 {—}{{---}}1 {–}{{--}}1,
}
\lstset{style=pythonkul}

% The complete-source appendix prints several thousand lines, so it uses a
% tighter variant of the same style: smaller type, no background tint (kinder
% to a printer over that many pages), and no frame, since the line numbers
% already mark the left edge. Everything else, including the line-breaking
% and the character mappings above, is inherited by restating it here.
\lstdefinestyle{sourcelisting}{%
    style=pythonkul,
    basicstyle=\ttfamily\scriptsize,
    backgroundcolor=\color{white},
    frame=none,
    xleftmargin=2.6em,
    numbersep=6pt,
    aboveskip=0.6\baselineskip,
    belowskip=0.4\baselineskip,
}

% Let cleveref name the new float types.
\crefname{algocf}{Algorithm}{Algorithms}
\Crefname{algocf}{Algorithm}{Algorithms}
\crefname{lstlisting}{Listing}{Listings}
\Crefname{lstlisting}{Listing}{Listings}

% ============================================================================
% "WHAT THE PROGRAM CHECKED" CALLOUT
% ============================================================================
% The recurring voice of the journey: a visually distinct box that separates
% narrative plus computational evidence from the formal theorem environments.
% It deliberately reuses the algorithmbox palette so the document feels unified.
\newtcolorbox{computernote}[1][Computer note]{%
    enhanced,
    breakable,
    colback=KULblauw3a!8,
    colframe=KULblauw3b!75,
    coltitle=white,
    fonttitle=\bfseries,
    title={#1},
    left=2mm, right=2mm, top=1mm, bottom=1mm,
    arc=2.5pt,
}

% ============================================================================
% CODE SIGNATURE CARDS  (the program's interface as the thesis's red wire)
% ============================================================================
% The computational architecture is the spine of this thesis, so its routines
% appear inline as readable, language-agnostic signature cards: a name, its
% inputs and output, a one-line "how", and an epistemic-role badge. No function
% bodies live in the running text; the full source accompanies the document.
\usepackage{tabularx}

% --- epistemic-role palette: the MODEL -> MEASURE -> PROVE/DISCOVER spine
%     spine, given one colour each so the role is legible at a glance. ----------
\definecolor{roleModel}{RGB}{82,189,236}      % KULblauw3a : the data model
\definecolor{roleMeasure}{RGB}{29,141,176}    % KULblauw1  : exact measurement
\definecolor{roleProve}{RGB}{60,160,80}       % green      : proved upper bounds
\definecolor{roleDiscover}{RGB}{220,140,40}   % orange     : discovered lowers
\definecolor{roleDriver}{RGB}{30,110,135}     % KULblauw3b : the single entry point

% A small filled badge; \colorbox keeps it robust inside tcolorbox titles.
\newcommand{\rb}[2]{{\setlength{\fboxsep}{1.6pt}%
    \colorbox{#1}{\textcolor{white}{\sffamily\bfseries\fontsize{7}{7}\selectfont #2}}}}
\newcommand{\rolebadge}[1]{\csname rolebadge@#1\endcsname}
\expandafter\newcommand\csname rolebadge@MODEL\endcsname{\rb{roleModel}{MODEL}}
\expandafter\newcommand\csname rolebadge@MEASURE\endcsname{\rb{roleMeasure}{MEASURE}}
\expandafter\newcommand\csname rolebadge@PROVE\endcsname{\rb{roleProve}{PROVE}}
\expandafter\newcommand\csname rolebadge@DISCOVER\endcsname{\rb{roleDiscover}{DISCOVER}}
\expandafter\newcommand\csname rolebadge@DRIVER\endcsname{\rb{roleDriver}{DRIVER}}

% Map each role to its accent colour so a card can tint its own frame, title bar,
% and left edge to match its badge. This makes the MODEL -> MEASURE -> PROVE /
% DISCOVER spine scannable down the page at a glance.
\expandafter\def\csname cardcol@MODEL\endcsname{roleModel}
\expandafter\def\csname cardcol@MEASURE\endcsname{roleMeasure}
\expandafter\def\csname cardcol@PROVE\endcsname{roleProve}
\expandafter\def\csname cardcol@DISCOVER\endcsname{roleDiscover}
\expandafter\def\csname cardcol@DRIVER\endcsname{roleDriver}

% ============================================================================
% PRIMER BOX: self-contained background on a standard tool, for a beginner.
% Visually distinct from the formal theorem environments and from computernote
% (which carries computational evidence): a soft green "learn this first" panel
% that an expert may skip, matching the appendix's "how to read this" ethos.
\newtcolorbox{primer}[1][Primer]{%
    enhanced, breakable,
    colback=roleProve!6,
    colframe=roleProve!55,
    coltitle=white,
    colbacktitle=roleProve!72,
    fonttitle=\bfseries,
    title={#1},
    left=2mm, right=2mm, top=1mm, bottom=1mm,
    arc=2.5pt,
}

% inline monospace code token
\newcommand{\code}[1]{{\ttfamily #1}}

% The optional first argument carries per-card colour overrides (set by
% \codecard from the role), the second is the card title.
\newtcolorbox{codecardbox}[2][]{%
    enhanced, breakable,
    colback=KULblauw1!3,
    colframe=KULblauw3b!55,
    boxrule=0.6pt, arc=1.5pt,
    left=2.5mm, right=2.5mm, top=0.8mm, bottom=0.8mm,
    colbacktitle=KULblauw3a!16,
    coltitle=KULblauw3b,
    fonttitle=\ttfamily\bfseries\small,
    title={#2},
    #1,
}
\newcommand{\cardlabel}[1]{\textcolor{KULblauw3b}{\sffamily\bfseries\scriptsize #1}}
% \codecard{signature}{ROLE}{input}{output}{how}
\newcommand{\codecard}[5]{%
\colorlet{cardaccent}{\csname cardcol@#2\endcsname}%
\begin{codecardbox}[%
    colframe=cardaccent!60,
    colbacktitle=cardaccent!16,
    coltitle=cardaccent!70!black,
    borderline west={2.2pt}{0pt}{cardaccent}]%
  {\strut #1\hfill\rolebadge{#2}}%
\setlength{\extrarowheight}{1.5pt}%
\begin{tabularx}{\linewidth}{@{}p{7mm}@{}X@{}}
\cardlabel{in}  & #3 \\
\cardlabel{out} & #4 \\
\cardlabel{how} & #5 \\
\end{tabularx}
\end{codecardbox}}

% --- the architecture spine, drawn once and referred back to throughout -------
\newcommand{\spinediagram}{%
\begin{tikzpicture}[>=Stealth,
  rolenode/.style={rounded corners=2.5pt, draw, very thick, align=center,
     minimum width=23mm, minimum height=12mm, font=\sffamily\footnotesize},
  flow/.style={->, thick, KULblauw3b!70}]
  \node[rolenode, fill=roleModel!16,    draw=roleModel]    (mod) at (0,0)     {\textbf{MODEL}\\[1pt]\scriptsize $\mu$ matrix};
  \node[rolenode, fill=roleMeasure!16,  draw=roleMeasure]  (mea) at (3.3,0)   {\textbf{MEASURE}\\[1pt]\scriptsize max-flow checker};
  \node[rolenode, fill=roleProve!18,    draw=roleProve]    (pro) at (7.0,0.85){\textbf{PROVE}\\[1pt]\scriptsize cut-counting, search};
  \node[rolenode, fill=roleDiscover!20, draw=roleDiscover] (dis) at (7.0,-0.85){\textbf{DISCOVER}\\[1pt]\scriptsize annealing, tabu};
  \draw[flow] (mod) -- (mea);
  \draw[flow] (mea) -- (pro);
  \draw[flow] (mea) -- (dis);
\end{tikzpicture}}

% --- the variant tree: root -> model (3) -> direction (2) -> separation (2) --
% Drawn twice, from the same coordinates, so the two figures are visibly one
% shape: plain in ch1 (a map before anything is proved), tinted by the
% trichotomy regime in ch4 (the same map, resolved). regimeProved/Conjectured/
% Open reuse the exact blue/red/amber the data figures already use for proved
% all-n / conjectured-with-certified-cases / open-lower-bound-only.
\colorlet{regimeProved}{KULblauw1}
\colorlet{regimeConjectured}{vertexred}
\colorlet{regimeOpen}{vertexorange}
\tikzset{
  vtroot/.style={rectangle, rounded corners=3pt, draw=KULblauw3b, very thick,
                 minimum width=26mm, minimum height=9mm, font=\sffamily\small\bfseries},
  vtmodel/.style={rectangle, rounded corners=3pt, draw=KULblauw3b, thick,
                  minimum width=20mm, minimum height=8mm, font=\sffamily\small},
  vtdir/.style={rectangle, rounded corners=2.5pt, draw=KULblauw3b!70,
                minimum width=19mm, minimum height=7mm, font=\scriptsize},
  vtleaf/.style={rectangle, rounded corners=2pt, draw=KULblauw3b!70,
                 minimum width=12mm, minimum height=6.5mm, font=\scriptsize},
  vtedge/.style={KULblauw3b!55, line width=0.9pt},
}
